\SourceLeaf{96}%
Porro habemus aequationes
\[
\begin{aligned}
(00)+(01)+(02)+(03)&=2n+1\\
(10)+(11)+(12)+(13)&=2n+1\\
(20)+(21)+(22)+(23)&=2n\\
(30)+(31)+(32)+(33)&=2n+1
\end{aligned}
\]
sive, adhibendo signa modo introducta, has tres (I):
\[
\begin{aligned}
h+i+k+l&=2n+1\\
2m+i+l&=2n+1\\
h+m&=n
\end{aligned}
\]
quarum itaque adiumento incognitas nostras iam ad duas reducere licet.

Aequationes reliquas e consideratione multitudinis solutionum congruentiae $1+\alpha+\beta+\gamma\equiv0$ derivabimus (per $\alpha,\beta,\gamma$, etiam hic indefinite numeros e complexibus $A,B,C$ resp. denotantes). Scilicet perpendendo \emph{primo}, $1+\alpha$ praebere $h,i,k,l$ numeros resp. ad $A,B,C,D$ pertinentes, et pro quovis valore dato ipsius $\alpha$ in his quatuor casibus resp. haberi solutiones $m,l,i,m$, multitudo omnium solutionum erit
\[
=hm+il+ik+lm
\]
\emph{Secundo} quum $1+\beta$ exhibeat $m,m,l,i$ numeros ad $A,B,C,D$ pertinentes, et pro quovis valore \emph{dato} ipsius $\beta$ in his quatuor casibus exstent solutiones $h,m,h,m$, multitudo omnium solutionum erit
\[
=hm+mm+hl+im
\]
unde derivamus aequationem
\[
0=mm+hl+im-il-ik-lm
\]
quae adiumento aequationis $k=2m-h$, ex (I) petitae, transit in hanc:
\[
0=mm+hl+hi-il-im-lm
\]
Iam ex aequationibus I habemus etiam $l+i=1+2h$, unde
\[
\begin{aligned}
2i&=1+2h+(i-l)\\
2l&=1+2h-(i-l)
\end{aligned}
\]

\SourceLeaf{97}%
Quibus valoribus in aequatione praecedente substitutis, prodit:
\[
0=4mm-4m-1-8hm+4hh+(i-l)^2
\]
Quodsi tandem pro $4m$ hic substituimus $2(h+m)-2(h-m)$ sive, propter aequationem ultimam in I, $2n-2(h-m)$, obtinemus:
\[
0=4(h-m)^2-2n+2(h-m)-1+(i-l)^2
\]
adeoque
\[
8n+5=(4(h-m)+1)^2+4(i-l)^2
\]
Statuendo itaque
\[
4(h-m)+1=a,\quad2i-2l=b
\]
fiet
\[
p=aa+bb
\]

Iam quum in hoc quoque casu $p$ unico tantum modo in duo quadrata, par alterum, alterum impar, discerpi possit, $aa$ et $bb$ erunt numeri prorsus determinati; manifesto enim $aa$ quadrato impari, $bb$ pari aequalis statui debet. Praeterea \emph{signum} ipsius $a$ ita erit stabiliendum, ut fiat $a\equiv1\ (\mathrm{mod.}\ 4)$, signumque ipsius $b$ ita, ut habeatur $b\equiv af\ (\mathrm{mod.}\ p)$, uti per ratiocinia iis, quibus in art. praec. usi sumus, prorsus similia facile demonstratur.

His praemissis quinque numeri $h,i,k,l,m$ per $a,b$ et $n$ ita determinantur:
\[
\begin{aligned}
8h&=4n+a-1\\
8i&=4n+a+2b+3\\
8k&=4n-3a+3\\
8l&=4n+a-2b+3\\
8m&=4n-a+1
\end{aligned}
\]
aut si expressiones per $p$ praeferimus, termini schematis $S$ per $16$ multiplicati ita se habebunt:
\[
\begin{array}{l|l|l|l}
p+2a-7&p+2a+4b+1&p-2a+1&p+2a-4b+1\\
p-2a-3&p-2a-3&p+2a-4b+1&p+2a+4b+1\\
p+2a-7&p-2a-3&p+2a-7&p-2a-3\\
p-2a-3&p+2a-4b+1&p+2a+4b+1&p-2a-3
\end{array}
\]

\SourceLeaf{98}%
\Article{21}
Postquam problema nostrum solvimus, ad disquisitionem principalem revertimur, determinationem completam complexus, ad quem numerus $2$ pertinet, iam aggressuri.

I. Quoties $p$ est formae $8n+1$, iam constat, numerum $2$ vel in complexu $A$ vel in complexu $C$ inveniri. In casu priori facile perspicitur, etiam numeros $\frac12(p-1)$, $\frac12(p+1)$ ad $A$ pertinere, in posteriori vero ad $C$. Iam perpendamus, si $\alpha$ et $\alpha+1$ sint numeri contigui complexus $A$, etiam $p-\alpha-1$, $p-\alpha$ tales numeros esse, sive, quod idem est, numeros complexus $A$ tales, quos sequatur numerus ex eodem complexu, binos semper associatos esse, ($\alpha$ et $p-1-\alpha$). Talium itaque numerorum multitudo, $(00)$, semper erit par, nisi quis exstat sibi ipse associatus, i. e. nisi $\frac12(p-1)$ ad $A$ pertinet, in quo casu multitudo illa impar erit. Hinc colligimus, $(00)$ imparem esse, quoties $2$ ad complexum $A$, parem vero, quoties $2$ ad $C$ pertineat. Sed habemus
\[
16(00)=aa+bb-6a-11
\]
sive statuendo $a=4q+1$, $b=4r$ (v. art. 14),
\[
(00)=qq-q+rr-1
\]
Quoniam igitur $qq-q$ manifesto semper par est, $(00)$ impar erit vel par, prout $r$ par est vel impar, adeoque $2$ vel ad $A$ vel ad $C$ pertinebit, prout $b$ est vel formae $8m$ vel formae $8m+4$. Quod est ipsum theorema, in art. 14 per inductionem inventum.

II. Sed etiam casum alterum, ubi $p$ est formae $8n+5$. aeque complete absolvere licet. Numerus $2$ hic vel ad $B$, vel ad $D$ pertinet, perspiciturque facile, in casu priori $\frac12(p-1)$ ad $B$, $\frac12(p+1)$ ad $D$, in casu posteriori autem $\frac12(p-1)$ ad $D$, $\frac12(p+1)$ ad $B$ pertinere. Iam perpendamus, si $\beta$ sit numerus ex $B$ talis, quem sequatur numerus ex $D$, fore etiam numerum $p-\beta-1$ ex $B$ atque $p-\beta$ ex $D$, i. e. numeros illius proprietatis binos associatos semper adesse. Erit itaque illorum multitudo, $(13)$, par, excepto casu, in quo unus eorum sibi ipse associatus est, i. e. ubi $\frac12(p-1)$ ad $B$, $\frac12(p+1)$ ad $D$ pertinet; tunc scilicet $(13)$ impar erit. Hinc colligimus, $(13)$ parem esse, quoties $2$ ad $D$, imparem vero, quoties $2$ ad $B$ pertineat. Sed habemus
\[
16(13)=aa+bb+2a+4b+1
\]

\SourceLeaf{99}%
sive statuendo $a=4q+1$, $b=4r+2$,
\[
(13)=qq+q+rr+2r+1
\]
Erit itaque $(13)$ impar, quoties $r$ par est; contra $(13)$ par erit, quoties $r$ est impar: unde colligimus, $2$ pertinere ad $B$, quoties $b$ sit formae $8m+2$, ad $D$ vero, quoties $b$ sit formae $8m+6$.

Summa harum investigationum ita enunciari potest:

\emph{Numerus $2$ pertinet ad complexum $A,B,C$ vel $D$, prout numerus $\frac12b$ est formae $4m$, $4m+1$, $4m+2$ vel $4m+3$.}

\Article{22}
In \emph{Disquisitionibus Arithmeticis} theoriam generalem divisionis circuli, atque solutionis aequationis $x^p-1=0$ explicavimus, interque alia docuimus, si $\mu$ sit divisor numeri $p-1$, functionem $\frac{x^p-1}{x-1}$ in $\mu$ factores ordinis $\frac{p-1}{\mu}$ resolvi posse adiumento aequationis auxiliaris ordinis $\mu$. Praeter theoriam generalem huius resolutionis simul casus speciales, ubi $\mu=2$ vel $\mu=3$, in illo opere art. 356--358 seorsim consideravimus, aequationemque auxiliarem a priori assignare docuimus, i. e. absque evolutione schematis residuorum minimorum potestatum alicuius radicis primitivae pro modulo $p$. Iam vel nobis non monentibus lectores attenti facile percipient nexum arctissimum casus proximi istius theoriae, puta pro $\mu=4$, cum investigationibus hic in artt. 15--20 explicatis, quarum adiumento ille quoque sine difficultate complete absolvi poterit. Sed hanc tractationem ad aliam occasionem nobis reservamus, ideoque etiam in commentatione praesente disquisitionem in forma pure arithmetica perficere maluimus, theoria aequationis $x^p-1=0$ nullo modo immixta. Contra coronidis loco adhuc quaedam alia theoremata nova pure arithmetica, cum argumento hactenus pertractato arctissime coniuncta, adiiciemus.

\Article{23}
Si potestas $(x^4+1)^{\frac12(p-1)}$ secundum theorema binomiale evolvitur, tres termini aderunt, in quibus exponens ipsius $x$ per $p-1$ divisibilis est, puta
\[
x^{2(p-1)},\quad Px^{p-1}\quad\text{atque }1
\]

\SourceLeaf{100}%
denotando per $P$ coëfficientem medium
\[
\frac{\frac12(p-1)\pdot\frac12(p-3)\pdot\frac12(p-5)\dotsfive\frac14(p+3)}{1\pdot2\pdot3\dotsfive\frac14(p-1)}
\]
Substituendo itaque pro $x$ deinceps numeros $1,2,3\dotsfour p-1$, obtinebimus per lemma art. 19
\[
\Sigma(x^4+1)^{\frac12(p-1)}\equiv-2-P
\]
At perpendendo ea quae in art. 19 exposuimus, insuperque, quod numeri complexuum $A,B,C,D$, ad potestatem exponentis $\frac12(p-1)$ evecti congrui sunt, secundum modulum $p$, numeris $+1,-1,+1,-1$ resp., facile intelligitur fieri
\[
\Sigma(x^4+1)^{\frac12(p-1)}\equiv4(00)-4(01)+4(02)-4(03)
\]
adeoque per schemata in fine artt. 18, 20 tradita
\[
\Sigma(x^4+1)^{\frac12(p-1)}\equiv-2a-2
\]
Comparatio horum duorum valorum suppeditat elegantissimum theorema: scilicet habemus
\[
P\equiv2a\ (\mathrm{mod.}\ p)
\]

Denotando quatuor producta
\[
\begin{array}{l}
1\pdot2\pdot3\mathpunct{.\,.\,.\,.\,.\,.\,.\,.\,.\,.\,.\,.\,.\,.\,.\,.\,.\,.\,.\,.\,.\,.\,.\,.\,.}\frac14(p-1)\\
\frac14(p+3)\pdot\frac14(p+7)\pdot\frac14(p+11)\dotsfive\frac12(p-1)\\
\frac12(p+1)\pdot\frac12(p+3)\pdot\frac12(p+5)\dotsfive\frac34(p-1)\\
\frac14(3p+1)\pdot\frac14(3p+5)\pdot\frac14(3p+9)\dotsfive(p-1)
\end{array}
\]
resp. per $q,r,s,t$, theorema praecedens ita exhibitur:
\[
2a\equiv\frac rq\ (\mathrm{mod.}\ p)
\]

Quum quilibet factorum ipsius $q$ complementum suum ad $p$ habeat in $t$, erit $q\equiv t\ (\mathrm{mod.}\ p)$, quoties multitudo factorum par est, i. e. quoties $p$ est formae $8m+1$, contra $q\equiv-t$, quoties multitudo factorum impar est, sive $p$ formae $8m+5$. Perinde in casu priori erit $r\equiv s$, in posteriori $r\equiv-s$. In utroque casu erit $qr\equiv st$, et quum constet, haberi $qrst\equiv-1$, erit $qqrr\equiv-1$,
